\documentclass[11pt]{article}

\usepackage[utf8]{inputenc}
\usepackage[T1]{fontenc}
\usepackage{lmodern}
\usepackage[authoryear]{natbib}
\usepackage{amsmath,amssymb}
\usepackage{booktabs}
\usepackage{array}
\usepackage{longtable}
\usepackage{geometry}
\usepackage{hyperref}
\usepackage{xurl}
\usepackage{microtype}
\usepackage{parskip}

\geometry{margin=1in}
\emergencystretch=2em

\hypersetup{
  pdftitle={Two cultures of mathematical writing},
  pdfauthor={Daniel Alami},
  pdfsubject={Research-move vocabularies, epistemic generation, and mechanization placement},
  colorlinks=true,
  linkcolor=blue,
  citecolor=blue,
  urlcolor=blue
}

\title{Two cultures of mathematical writing\\
\large A corpus study of research-move vocabularies}
\author{Daniel Alami}
\date{June 2026}

\begin{document}
\maketitle

\begin{abstract}
Research systems need to propose useful next moves before final claims exist.
We treat such epistemic generation as mechanization placement. We name and
stress-test recurring research moves, then assign each one to automation, an
evidence contract, a deterministic check, a policy feature, or human judgment
only when the local contract warrants that placement. Our evidence comes from
one operated research corpus and three connected empirical layers. A
corpus-level operationalization of Gowers's theory-builder/problem-solver
distinction finds a symmetric own-corpus advantage: a theory-builder
vocabulary covers theory-builder arcs at 58.1\% and problem-solver arcs at
20.7\%, while a problem-solver vocabulary covers problem-solver arcs at 65.9\%
and theory-builder arcs at 19.1\%. A same-culture random split produces only a
3.0 percentage-point gap. A later eight-subfield remine then revises the binary
into a broader structural vocabulary: six shared-core operations, eight broadly
shared operations, and four peripheral operations over 1,214 moves. That
vocabulary transfers beyond its original setting and exposes measurable
reliability structure, with a methodological warning emerging from a PDE
rescore. Prompt and workflow experiments locate the interface that makes the
vocabulary operational. Passive labels and broad menus carry little causal
force, while typed evidence contracts, source-bound residual-to-check edges,
rejected confusers, and checked action schemas change downstream behavior in
replay and synthetic workflow settings. Meta-language experiments add a further
result: meta-operator semantics are separable from adjacent lower-tier moves and
can sharpen already-parsed evidence artifacts. A research-move vocabulary
becomes scientifically useful when it assigns a move to the right
role, recognition, human review, evidence contract, deterministic check,
policy feature, action schema, or judgment.
\end{abstract}

\section{Introduction}

To verify a claim is to ask whether it is adequately supported. But before any
claim exists, a research process has to pick its next move: reformulate the
problem, try a special case, draw an analogy, look for a falsifier, introduce an
auxiliary object, decompose, mark an evidence boundary, or stop a branch.

Two curated mathematical styles turn out to favor different research-move
vocabularies, and the difference survives cross-corpus scoring against a
same-culture negative control. Gowers's theory-builder/problem-solver distinction
\citep{gowers2000} thus gets an empirical, operational form, without forcing all
of mathematics into two types.

We then widen the binary into a structural vocabulary of research moves.
Theory-building and problem-solving mark a real but incomplete split: they stress
different bands inside a larger space of reformulation, abstraction,
decomposition, local-to-global assembly, invariance, translation, approximation,
and constraint propagation.

This vocabulary earns its keep at one specific point in a research system. Passive
labels and long menus do little on their own, typed evidence fields,
residual-to-check edges, rejected nearest confusers, and executable action schemas
do the work. Research-move language turns operational only when it changes what
the system checks, blocks, turns into an artifact, or executes next.

Generation, then, is a sequence of research moves, judged by whether the system
picked, checked, and executed the right one. A generic evaluation vocabulary can
miss domain-specific reasoning that typed intermediate representations recover.

\section{Evidence base, reproducibility, and method}

Our unit of analysis is the research move: a bounded step that changes a problem
representation, admissible object, proof strategy, evidence standard,
decomposition, validation frame, or next action. A vocabulary covers a move when
a frozen operation fits it mechanistically under our scoring rubric.

Coverage is only a diagnostic. High coverage can index transfer, but it can just
as easily signal force-fitting, so we optimize nothing against it directly. Every
coverage number travels with cross-corpus scoring, negative controls, adversarial
raters, decoys, wrong-card controls, source-cluster blocking, and downstream
consequence checks. Raw fit to a vocabulary is weak evidence, fit that survives
the wrong-card, decoy, source-only, and downstream-consequence checks is what
carries weight.

The argument draws on several tiers of evidence, each with a different
role. Saved-output corpus results carry the main quantitative claims: the
two-cultures split, the negative control, structural-vocabulary coverage,
and out-of-distribution transfer. Cross-family and multi-rater audits
supply the reliability checks, especially for the PDE correction and the
catalog-confusion analysis. Replay and synthetic-workflow experiments give
the mechanism evidence for evidence carriers, residual-to-check edges,
boundary cards, and action schemas. Human and production-trace validation
is the next layer, for expert use and prospective deployment.


Durable records back the core percentages: stored move lists, frozen
vocabularies, scoring files, JSON verdicts, and reproducer scripts. Later
workflow experiments serve as mechanism evidence, showing which carrier changes
downstream behavior under controlled conditions. The reproducibility packet
bundles the saved-output verifiers for the headline corpus claims, the PDE
joint-vocabulary scoring artifacts, the meta-language score files, and the
compact workflow replay scorers. Each main quantitative claim rests on a
recomputable artifact, so the record stands on its own without the underlying
model transcripts.

\section{A corpus-level two-cultures result}

Theory-building arcs make up our first corpus: Wiles, Grothendieck, Lurie,
Scholze, Riemann, Newton, Einstein, and Polya/Hadamard. The theory-builder
vocabulary we mine from them includes foundational object redefinition,
cross-domain unification, parameter-space internalization, vocabulary lifting,
strategic specialization, diagonal self-application, and concept revision
\citep{polya1945,lakatos1976}.

We collect problem-solving arcs in a sister corpus: Erdos discrepancy,
Green--Tao, Hales--Jewett Polymath, Szemeredi regularity, Roth, Behrend,
Furstenberg, and Ramsey/Erdos--Szekeres. Its problem-solver vocabulary
emphasizes partitioning, governed iterative refinement, theorem-use transfer,
rank induction, and estimate chaining.

\begin{table}[ht]
\centering
\caption{Cross-distribution coverage for the two mined vocabularies.}
\label{tab:two-cultures}
\begin{tabular}{lrr}
\toprule
Vocabulary & Theory-builder corpus & Problem-solver corpus \\
\midrule
Theory-builder vocabulary & 58.1\% & 20.7\% \\
Problem-solver vocabulary & 19.1\% & 65.9\% \\
\bottomrule
\end{tabular}
\end{table}

Each vocabulary fits its own corpus far better than the other, by 42.1
percentage points on average, and the advantage runs symmetrically in both
directions. Split one culture at random and the gap drops to 3.0 points, far
below the curated two-cultures gap. So two mathematical styles favor different
research-move vocabularies that a random partition of one style cannot reproduce,
and the Gowers distinction picks up an empirical, operational form. This
operationalizes the axis descriptively, it does not pin it down causally against
every corpus-construction factor, since subfield, era, and selection effects may
all feed the gap. The narrower claim is the one we make: the curated contrast
yields a measurable research-move difference that a same-style random split does
not.

Studies of mathematical practice, proof reading, examples, collaboration,
consensus, subfield prestige, %
\citep{cranshaw2011polymath,gargiulo2022polymath,weber2014proofreading,lockwood2016examples,mejia2021explanatory,wagner2022consensus,schlenker2020prestige}
stop short of a corpus-level operationalization of Gowers's
theory-builder/problem-solver distinction. None of that adjacent work tests
directly whether the two styles differ by cross-coverage of research-move
vocabularies, and that is the test the result above supplies.

\section{From two cultures to a structural vocabulary}

The two-cultures result is only a starting point. A follow-up eight-subfield
analysis covered 1,214 moves across 64 mathematical arcs and compressed them into
six shared-core operations, eight broadly shared operations, and four peripheral
operations, a taxonomy already wider than the binary.

\begin{table}[ht]
\centering
\caption{Layered structural vocabulary from the eight-subfield remine.}
\label{tab:structural-vocabulary}
\begin{tabular}{p{0.18\linewidth}p{0.74\linewidth}}
\toprule
Tier & Operations \\
\midrule
Shared core & Problem Reformulation and Reduction, Generalization and Abstraction, Decomposition and Recomposition, Local-to-Global Assembly, Canonical Form and Invariance, Cross-Domain Translation. \\
Broadly shared & Iterative Refinement, Recursive Decomposition, Duality and Adversarial Framing, Layered Approximation and Convergence, Extremal Method, Probabilistic and Stochastic Methods, Dimensional and Structural Lifting, Constraint Imposition and Propagation. \\
Peripheral & Characterization by Obstruction, Internalization and Self-Reference, Axiomatization and Foundational Repair, Controlled Universe Extension. \\
\bottomrule
\end{tabular}
\end{table}

This revision generalizes the original result: theory-builder and problem-solver
mathematics stress different bands inside a larger structural language. Some moves
are shared, others are subfield idioms, and several seemingly culture-specific
moves turn out to be aliases once the remine widens.

Transfer also reaches past the original mathematical corpus. A business held-out
set reaches 57.9\% shared-plus-broad coverage, and four sparse 2026 specialist
papers reach 75.2\%. We record scope, residuals, and confusers to bound this
portability as measured, keeping the coverage well short of any claim to
universality.

\section{Methodological warning from the PDE rescore}

The PDE rescore carries a methodological lesson for evaluating frontier agents
and research systems: domain-specific reasoning can look absent whenever an
evaluation vocabulary omits the operations that carry it. A 12.5\% adversarial
score on a post-cutoff PDE paper first looked like a sharp capability boundary for
the structural vocabulary. Auditing it exposed a different mechanism, the score
came from a single vocabulary, a single adversarial rater, and one hard PDE paper.
Rescore against the joint structural vocabulary plus the PDE estimate-craft
vocabulary, and the same paper draws 95.8--100\% coverage of its moves from four
raters across two model families under the covered-vs-none metric. Across the
five-paper corpus, two-family coverage reached 116/118 moves, or 98.3\%, with a
0\% decoy match rate.


Test frontier agents on domain-specific reasoning using only generic reasoning
categories, and you can manufacture artificial capability boundaries. The
expanded language plus PDE estimate-craft receipts covered this former stress
case under multi-rater, decoy-controlled scoring. Throughout, we report
covered-vs-none coverage and operation-identity reliability, and keep strict
full-only coverage as a calibration-sensitive secondary metric.

The rescore cuts both ways, and the sharper objection runs the other
direction: if the wrong vocabulary drives a fixed artifact to 12.5\% and
the right one to nearly 100\%, coverage is measuring the vocabulary, not
the paper. On a single artifact that is true, which is why no claim here
rests on one paper's coverage. The weight falls instead on three controls
that vocabulary choice cannot satisfy at once. A same-culture negative
control blocks easy inflation: a generous vocabulary applied to two papers
from the same culture does not reproduce the cross-culture split. Held-out
transfer blocks overfitting: coverage holds on papers the vocabulary was
never tuned against, including business and 2026 specialist corpora. The
0\% decoy-match rate blocks looseness: the vocabulary does not light up on
planted moves it should miss. A vocabulary tuned to inflate one paper's
score breaks at least one of these three, so the coverage numbers count
only where all three hold.

\section{Carrier selection in agent-facing tests}

Knowing the vocabulary is one thing, getting an agent to act on it is another. We
supplied research-move vocabulary three ways, as prompt text, as feature prose, and
as a broad catalogue, and none of them moved the tested outcomes much.

On an exact-answer benchmark, every prompt arm (baseline, static catalogue,
one-primitive, diagnose-then-solve) scored the same 11/26, and a placebo 8/26: no
lift from the prompt text. On blind route-choice probes, generic reasoning won 7 of
10 comparisons while each vocabulary-supplied arm won 6 of 10, so the vocabulary
changed how the agent explained its choice more than the choice itself. A
specialized card for a move beat its schema and swap controls in 0 of 24 cases,
routing the choices to execution did not rescue the vocabulary-supplied routes (51
votes to ordinary routes, 47 to schema routes, 28 to vocabulary routes), and a
consequence-ranking test gave the card only a 4.25/5 to 4.04/5 edge, ahead of the
strongest control in just 2 of 20 source clusters. Supplied this way, the
vocabulary is close to inert.

It begins to matter only once the correct move has already been chosen. An
action-oriented catalogue merely matched source-only reasoning (6.50/10 each,
against 6.00/10 for a descriptive summary), while a correct action card scored
8.875/10 against 4.375/10 for a plausible wrong one. So the open problem is not
naming more moves, it is routed selection: picking the right move, rejecting the
nearest confuser, and attaching the evidence that makes the choice executable.

\section{Implications for research-agent evaluation}

The experiments point to three design requirements for evaluating research agents
on open-ended scientific or mathematical work.

Final-answer accuracy is too coarse for epistemic generation. A research agent can
fail by picking the wrong object, paying the wrong evidence debt, over-updating
from a narrow receipt, dropping a necessary falsifier, or pressing on after a
branch should have stopped. A single final-answer score routinely hides such
failures.

An evaluation vocabulary has to include the domain operations that carry the
reasoning. In the PDE rescore, a generic structural vocabulary understated
capability the moment estimate-craft moves fell out of the scoring language. The
same point reaches past PDE to clinical, legal, institutional, and workflow
settings, each of which carries admissibility and authority conditions that
generic reasoning labels can miss.

Useful evaluation artifacts preserve action-relevant state, so the most
informative records here are typed fields: selected residual edge, rejected
nearest confuser, source evidence, satisfied and unsatisfied evidence
obligations, action schema, step index, invariant check, and later outcome. These
fields show whether a system executed the right research move where a plausible
description of it would have passed an accuracy score.

\section{Meta-language}

In one track we ask whether some moves operate a level above the base structural
vocabulary: changing the admissible criteria, making an equivalence class primary,
promoting an auxiliary object into the main object of study, or reframing the
question itself.

Across these families one pattern holds: meta cards are separable from
adjacent lower-tier forms, typed parsing identifies the family, and a
target card sharpens an already-parsed artifact, lifting frozen
typed-frame artifacts from 6.91/8 to 7.91/8. Easy cases reach ceiling
without testing downstream use, and content matters more than packaging.

Compact meta-level semantics improve already-parsed evidence artifacts and tell
frame-changing moves apart from adjacent lower-tier ones. Wrong meta cards can
mis-shape downstream artifacts, which gives the meta layer causal content. The
evidence supports three candidate meta families and argues against expanding to a
fourth before the data warrant it.

\section{From vocabulary to action schemas}

The strongest later evidence arrives when the vocabulary is compiled into an
action schema the agent executes, with the label form held back as a control.

Boundary-card experiments show the pattern. Agents act correctly from a fully
specified boundary card that states the satisfied evidence obligation,
unsatisfied evidence obligation, permitted update, blocked update, next-action
rule, and false-reading confuser. Letting the model extract that card raw is
unreliable, model-only validation improves but still waves through unsafe cards,
rule-backed source-cue validation plus action-schema fields recover correct
behavior on the tested packets.

Process-control experiments show the same architecture. A compiled action
schema, selected residual edge, rejected nearest confuser, source evidence, step
index, required next action, reached 1.0 action accuracy on the held-out
two-step packet. Label-only controls failed, and so did free-form automatic
compilation. Typed class selection helped but stayed unsafe until we added
source-cue checks and deterministic lowering. Open-set refusal kept outside cases
from being forced into known classes.

The resulting architecture is a research-agent contract:
\[
\begin{aligned}
\text{source facts} &\rightarrow \text{residual/evidence carrier}
\rightarrow \text{nearest confuser} \\
&\rightarrow \text{action program}
\rightarrow \text{deterministic check}
\rightarrow \text{outcome trace}.
\end{aligned}
\]
This sits next to algorithm selection, metareasoning, selective prediction,
mixture-of-experts routing, prompt-programming systems, and workflow-agent
evaluation \citep{rice1976,xu2008satzilla,jacobs1991adaptive,russell1991rational,geifman2017selective,chen2023frugalgpt,ong2024routellm,hu2024routerbench,flowagent2025,sopbench2025,agentcompany2024}.
As a runtime contract for research agents, it binds local source facts to the
next inspectable action.

\section{Mechanization placement}

Across the math corpus, structural vocabulary, agent-prompt tests, meta-language
probes, and action-schema experiments, one result recurs: placement. A single
research move can sit in several different roles.

\begin{table}[ht]
\centering
\caption{Placement roles for research-move vocabulary.}
\label{tab:placement}
\begin{tabular}{p{0.21\linewidth}p{0.36\linewidth}p{0.33\linewidth}}
\toprule
Placement & Appropriate use & Failure mode \\
\midrule
Recognition language & Naming and comparing research moves. & Mistaken for a recipe. \\
Human review aid & Helping reviewers remember obligations and confusers. & Measured by self-reported usefulness instead of blinded decisions. \\
Evidence contract & Forcing an output into inspectable fields. & Shape is imitated without source alignment. \\
Deterministic check & Enforcing a crisp pass/fail condition. & Check overreaches beyond its contract. \\
Policy feature & Supplying state to a decision rule. & Retrospective labels mistaken for prospective policy quality. \\
Action schema & Binding source facts to next actions and stop rules. & Unsafe if parsed free-form or without refusal. \\
Judgment & Handling semantic frontier choices. & Hidden behind premature automation. \\
\bottomrule
\end{tabular}
\end{table}

Place each move at the strongest level its evidence supports. Structural
vocabulary is the recognition layer, typed receipts and nearest-confuser fields
form the evidence-contract layer, deterministic gates earn their place once the
contract is crisp. Human judgment stays necessary wherever a move is semantic,
frontier-dependent, or mathematically substantive.

\section{Future validation}

Five bounded claims now stand on the evidence: the two-cultures
operationalization holds under its negative control, the structural vocabulary
ports across fields, the PDE audit corrects the earlier boundary, the
meta-language artifact result is scoped to typed conditions, and carrier
selection accounts for where mechanization pays off. Further experiments would
strengthen external validity and prospective-use claims.

Four follow-ons carry the most value: (i) a blinded human or expert audit of
structural-vocabulary assignments and checklist usefulness, (ii) production or
historical-trace replay testing whether compiled contract fields reduce
downstream cost, missed obligations, false stops, or wrong actions relative to
strong context baselines, (iii) a clean meta-language downstream-consequence
test with generic headroom, and (iv) a fresh blind PDE or broader scientific
corpus audit using covered-vs-none coverage, operation-identity reliability,
decoys, and human review.

\section{Broader impact}

Measuring intermediate research moves, evidence payments, confuser rejection, and
action schemas sharpens evaluation discipline for research agents. A reviewer can
then see where a final-answer score or a broad reasoning label would have papered
over a domain-specific limitation, and the claims one can make about agent
capability tighten accordingly.

A structural label can also be misused as an authority signal, dressing up a weak
research step to look more systematic than it is. Placement already supplies the
mitigation: a label stays advisory until it ties to source evidence, falsifiers,
deterministic checks, or downstream outcomes. Domain-specific evaluation should
report its residuals and confusers, leaving every move that resists the nearest
available category visibly uncategorized.

\section{Limitations}

We cover one research system and curated corpora, which gives us a measured
vocabulary and a mechanization account for research moves.

Most of our coverage judgments are model-mediated, so we shore up several claims
with cross-family and multi-rater audits. Independent human and expert audits are
the natural next validation layer, here, human audit remains a future
external-validity test whose evidentiary value is still to be collected.

Several agent-facing experiments are small-N mechanism tests. Their point
estimates diagnose mechanism, so reading them as population estimates would
overreach, and small differences should not carry inferential weight. The main
weight falls on large contrasts, negative controls, decoys, wrong-card tests, and
cases where a carrier change alters the downstream artifact.

Coverage measures structural recognition, so generative competence, frontier
selection, and mathematical correctness need separate tests.

The evidence supports a compact core vocabulary plus field-specific receipts and
confuser guards, while also surfacing overlap, low-use operations, and
calibration-sensitive full/partial distinctions in some catalog versions.

For meta-language, the evidence supports separability and artifact refinement
under typed conditions, with incremental routing and downstream consequence
tests as the next extension.

Action-schema positives are strongest in replay, synthetic, and rule-checked
settings, with production-trace tests as the next external-validity step.

\section{Conclusion}

A Gowers-style question opened the inquiry: do theory-building and
problem-solving arcs favor different research-move vocabularies? In this corpus,
they do. The binary sits inside a structural vocabulary whose usefulness depends
on placement.

Passive labels have little causal force, a correct compact schema can help once
the move is selected, meta-language sharpens typed artifacts, checked action
schemas change downstream behavior. So the discipline is to name recurring
research moves, test their boundaries, and mechanize the parts whose contracts can
be made inspectable.

\appendix
\section{Evidence map}

\begingroup
\scriptsize
\begin{longtable}{p{0.30\linewidth}p{0.60\linewidth}}
\caption{Local artifacts that back each claim family.}
\label{tab:evidence-map}\\
\toprule
Claim family & Saved artifact or source \\
\midrule
\endfirsthead
\toprule
Claim family & Saved artifact or source \\
\midrule
\endhead
Two-cultures cross-coverage & \texttt{evidence/gp216\_queries/gp216\_cross\_distribution\_full.json} \\
Headline verifier & \texttt{evidence/reproducers/verify\_gp216\_claims.py} \\
Same-culture negative control & \texttt{evidence/gp216\_queries/gp216\_negative\_control.json} \\
Four-operation compression & \texttt{evidence/gp216\_queries/gp216\_4op\_compression.json} \\
Eight-subfield structural vocabulary & \texttt{evidence/gp216\_queries/gp216\_pathA\_8sf\_cluster.json} \\
Business held-out transfer & \texttt{evidence/gp216\_queries/gp216\_business\_held\_out\_OOD.json} \\
Sparse 2026 specialist transfer & \texttt{evidence/gp216\_queries/gp216\_sparse\_coverage\_OOD.json} \\
PDE correction & \texttt{experiments/joint\_vocab\_corpus\_ood\_vC01\_20260521/} \\
Catalog reliability stress test & \texttt{experiments/v128b\_catalog\_stats\_20260521/} \\
Catalogue surface tests & \texttt{experiments/source\_packet\_catalog\_surface\_v14/} and follow-up catalogue/card tests in \texttt{research\_log.md} \\
Meta-language tests & Meta-card, typed-frame, and packaging-ablation tests in \texttt{research\_log.md} \\
Boundary-card and action-schema tests & Boundary-card and process-control tests in \texttt{research\_log.md} \\
Pattern/action contract tests & Pattern/action contract synthesis in \texttt{research\_log.md} and \texttt{research\_roadmap.md} \\
Literature audit & \texttt{literature\_audit.md} \\
Claim guardrail & \texttt{claim\_evidence\_matrix.md} \\
\bottomrule
\end{longtable}
\endgroup

\bibliographystyle{plainnat}
\bibliography{refs}

\end{document}
